\appendix
\section*{Appendices}
\addcontentsline{toc}{section}{Appendices}
\renewcommand{\thesubsection}{\Alph{subsection}}
\setcounter{subsection}{0}

\subsection{An integral identity for the hypergeometric function}
\label{app:hyp}

The route to \cref{eq:Psidef} taken in \cref{sec:absbirthdeath} expands the
characteristic $x(\tau)$ in a geometric series and integrates term by term. An
alternative is to integrate the quotient \cref{xabd} directly, splitting the
numerator into two pieces, each of the form
\begin{equation}
\label{eq:Iform}
I(\tau) = \int\frac{\ee^{h\tau}}{p+q\,\ee^{k\tau}}\,\dd\tau ,
\end{equation}
with $h,k,p,q$ constant. That route is longer and easier to get wrong --- the two
pieces carry different prefactors, and dropping one produces a formula that still
satisfies the initial condition while violating $G(1,1,t)=1$ --- but the identity it
rests on is useful in its own right, so we record and prove it here.

\begin{proposition}
\label{prop:hypIden}
For $h/k>0$ and $-\tfrac{q}{p}\ee^{k\tau}\notin[1,\infty)$,
\begin{equation}
\label{hypIden}
\int\frac{\ee^{h\tau}}{p+q\ee^{k\tau}}\,\dd\tau
= \frac{\ee^{h\tau}}{ph}\,{}_2F_1\lrb{1,\frac{h}{k};\frac{h}{k}+1;\,-\frac{q}{p}\ee^{k\tau}} + \text{const}.
\end{equation}
\end{proposition}

\begin{proof}
We use the Euler integral representation of the hypergeometric function,
\begin{equation}
\label{idenHyp}
{}_2F_1(a,b;c;z) = \frac{\Gamma(c)}{\Gamma(b)\Gamma(c-b)}\int^1_0 u^{b-1}(1-u)^{c-b-1}(1-zu)^{-a}\,\dd u ,
\end{equation}
$\Gamma$ being the standard gamma function.

Start from \cref{eq:Iform} and substitute $v = \ee^{k\tau}$, so that
$\dd\tau = k^{-1}v^{-1}\dd v$:
\begin{align*}
I &= \frac{1}{k}\int\frac{v^{h/k}}{p+qv}\,v^{-1}\dd v
= \frac{1}{pk}\int\frac{v^{h/k-1}}{1+\frac{q}{p}v}\,\dd v .
\end{align*}
Fix the antiderivative by writing this as a definite integral from $0$ to $v$ in a
dummy variable $w$ (the integral converges at the lower limit precisely because
$h/k>0$):
\begin{equation*}
I = \frac{1}{pk}\int^v_0\frac{w^{h/k-1}}{1+\frac{q}{p}w}\,\dd w .
\end{equation*}
Now rescale, $w = vu$, so that $\dd w = v\,\dd u$ with $v$ held fixed:
\begin{align}
I &= \frac{1}{pk}\int^1_0\frac{(vu)^{h/k-1}}{1+\frac{q}{p}vu}\,v\,\dd u
= \frac{v^{h/k}}{pk}\int^1_0\frac{u^{h/k-1}}{1+\frac{q}{p}vu}\,\dd u . \label{backat}
\end{align}
The remaining integrand matches that of \cref{idenHyp},
\begin{equation*}
\int^1_0 u^{h/k-1}(1-u)^{0}\Bigl(1-\bigl(-\tfrac{q}{p}v\bigr)u\Bigr)^{-1}\dd u
= \int^1_0 u^{b-1}(1-u)^{c-b-1}(1-zu)^{-a}\,\dd u ,
\end{equation*}
with $b = h/k$, $c = h/k+1$, $a = 1$ and $z = -\tfrac{q}{p}v$. Hence
\begin{align*}
\int^1_0\frac{u^{h/k-1}}{1+\frac{q}{p}vu}\,\dd u
&= \frac{\Gamma\!\lrb{\frac{h}{k}}\Gamma(1)}{\Gamma\!\lrb{\frac{h}{k}+1}}\,{}_2F_1\lrb{1,\frac{h}{k};\frac{h}{k}+1;-\frac{q}{p}v}
= \frac{k}{h}\,{}_2F_1\lrb{1,\frac{h}{k};\frac{h}{k}+1;-\frac{q}{p}v},
\end{align*}
using $\Gamma(z+1) = z\Gamma(z)$. Substituting into \cref{backat},
\begin{equation*}
I = \frac{v^{h/k}}{ph}\,{}_2F_1\lrb{1,\frac{h}{k};\frac{h}{k}+1;-\frac{q}{p}v},
\end{equation*}
and restoring $v = \ee^{k\tau}$ gives \cref{hypIden}.
\end{proof}

\begin{remark}
Applied to the two pieces of \cref{xabd} --- the first with $h=\alpha$, the second
with $h = \alpha+b_1$, and both with $k=b_1$, $p=(s-x_-)$, $q=-(s-x_+)$ ---
\cref{prop:hypIden} reproduces \cref{eq:Psihyp} after the prefactors
$\alpha x_+(s-x_-)$ and $-\alpha x_-(s-x_+)$ are carried through correctly. The
factor $(s-x_+)/(s-x_-)$ attached to the second piece is easily lost, and its loss is
invisible at $t=0$; checking $G(1,1,t)=1$ catches it immediately.
\end{remark}

\subsection{Extracting state probabilities from a generating function}
\label{app:extract}

A generating function defined by a formula rather than as an explicit series --- for
example
\begin{equation}
\label{examplePgf}
G(t;x,y) = \Bigl(\bigl(1-\ee^{-\alpha t}\bigr)x + \ee^{-\alpha t}y\Bigr)^n
\end{equation}
--- can always be converted back to the form
\begin{equation}
\label{series1}
G(t;x,y) = \sum^{\infty}_{i=0}\sum^{\infty}_{j=0} p_{i,j}(t)\,x^iy^j ,
\end{equation}
recovering the time-dependent state probabilities. All that is required is that
$G$ be repeatedly partially differentiable in $x$ and $y$.

The reason is the bivariate Maclaurin expansion,
\begin{equation*}
f(x,y) = \sum^{\infty}_{k=0}\frac{1}{k!}\sum^{k}_{\rho=0}\binom{k}{\rho}\,\partial_x^{\rho}\partial_y^{\,k-\rho}f\Big|_{(0,0)}\;x^{\rho}y^{k-\rho} ,
\end{equation*}
whose coefficient of $x^iy^j$ (taking $k=i+j$, $\rho=i$) is
$\frac{1}{(i+j)!}\binom{i+j}{i}\partial_x^i\partial_y^{\,j}f\big|_{(0,0)}$. Hence
\begin{equation}
\label{stateProb}
p_{i,j}(t) = \frac{1}{i!\,j!}\,\partial_x^i\partial_y^{\,j}G\Big|_{(0,0)} .
\end{equation}

\subsubsection*{The multilinear case}

Both \cref{eq:Gabsonly} and \cref{eq:Gabsdeath} have the form
\begin{equation}
\label{eq:multilinear}
G(t;x,y) = \bigl(c(t) + f(t)\,x + g(t)\,y\bigr)^n ,
\end{equation}
for which the derivatives are immediate. Each differentiation lowers the power by
one and brings down a factor, so
\begin{equation}
\label{eq:derivs}
\partial_x^{a}\partial_y^{\,b}G = \frac{n!}{(n-a-b)!}\,f^{a}g^{b}\bigl(c+fx+gy\bigr)^{n-a-b}
\qquad (a+b\le n),
\end{equation}
and zero for $a+b>n$. Evaluating at the origin and applying \cref{stateProb},
\begin{equation}
\label{stateProbSpec}
p_{i,j}(t) = \frac{n!}{i!\,j!\,(n-i-j)!}\;f(t)^{i}\,g(t)^{j}\,c(t)^{\,n-i-j},
\qquad i+j\le n,
\end{equation}
and $p_{i,j}=0$ otherwise. This is the trinomial distribution, which is exactly what
it should be: the $n$ particles are independent, and at time $t$ each is in one of
three states with probabilities $c$ (neither inside nor outside --- that is, dead),
$f$ (inside) and $g$ (outside). Note in particular that the falling factorial
$n!/(n-i-j)!$ in \cref{eq:derivs} is \emph{not} the same as $n^{i+j}$; the two agree
only to leading order in $n$, and using the latter breaks normalisation.

\subsubsection*{Applied to the absorption models}

For absorption only, \cref{examplePgf}, we have $c=0$, $f = 1-\ee^{-\alpha t}$ and
$g = \ee^{-\alpha t}$. The factor $c^{\,n-i-j}$ then vanishes unless $i+j=n$, and
\cref{stateProbSpec} gives
\begin{equation*}
p_{i,j}(t) = \binom{n}{i}\bigl(1-\ee^{-\alpha t}\bigr)^{i}\bigl(\ee^{-\alpha t}\bigr)^{j},
\qquad i+j = n,
\end{equation*}
so $\xt\sim\text{Binomial}\bigl(n,\,1-\ee^{-\alpha t}\bigr)$ and $\yt = n-\xt$, as
anticipated in \cref{sec:absonly}.

For absorption--death, \cref{eq:Gabsdeath}, the three functions are
\begin{equation*}
c(t) = 1-\frac{\mu}{\mu-\alpha}\ee^{-\alpha t} + \frac{\alpha}{\mu-\alpha}\ee^{-\mu t},
\qquad
f(t) = \frac{\alpha}{\mu-\alpha}\bigl(\ee^{-\alpha t}-\ee^{-\mu t}\bigr),
\qquad
g(t) = \ee^{-\alpha t},
\end{equation*}
which are precisely $p_{0,0}$, $p_{1,0}$ and $p_{0,1}$ of \cref{eq:absdeathstates},
and
\begin{equation*}
p_{i,j}(t) = \frac{n!}{i!\,j!\,(n-i-j)!}\;f(t)^{i}\,g(t)^{j}\,c(t)^{\,n-i-j}.
\end{equation*}
Each particle independently ends up dead, absorbed, or still outside, with those
three probabilities; the joint law is trinomial. That the generating function
factorises in this way is of course just \cref{indRelPGF} read backwards.

The absorption--birth--death generating function of \cref{prop:Gabd} is not of the
form \cref{eq:multilinear}, and there \cref{stateProb} must be applied to
\cref{eq:Gabd} directly, or the coefficients extracted numerically by Cauchy
integrals on circles $|x|=|y|=\rho<1$.
